\clearpage
\begin{summarybox}
\textbf{\textcolor{CSummary}{一句话核心}}
\textbf{先会写切线，再把“给定点”和“切点”的坐标交换：这就是极线公式。}
切点弦是它在\textbf{有两个不同实切点}时的几何表现，并非每条极线都能看成切点弦。

\medskip
\textbf{\textcolor{CSummary}{三种常用方程}}
\begin{align*}
 \text{椭圆：}&\quad\boxed{\frac{x_0x}{a^2}+\frac{y_0y}{b^2}=1},\\
 \text{双曲线：}&\quad\boxed{\frac{x_0x}{a^2}-\frac{y_0y}{b^2}=1},\\
 \text{抛物线 }y^2=2px\text{：}&\quad\boxed{y_0y=p(x+x_0)}.
\end{align*}
一般式 $Ax^2+Bxy+Cy^2+Dx+Ey+F=0$：
\[
 \boxed{\begin{aligned}
 &Ax_0x+\frac{B}{2}(x_0y+y_0x)+Cy_0y\\
 &\qquad+\frac{D}{2}(x+x_0)+\frac{E}{2}(y+y_0)+F=0.
 \end{aligned}}
\]

\medskip
\textbf{\textcolor{CSummary}{四法如何选择}}
\begin{itemize}[leftmargin=*,itemsep=0.3em]
 \item \textbf{最低门槛：}先从圆的“半径垂直切线”理解；椭圆靠伸缩，双曲线、抛物线靠代回去检验。
 \item \textbf{考场最快：}已知切线公式，令给定点 $P$ 在切线上，交换两组坐标。
 \item \textbf{证明步骤最少：}设直线与曲线联立得 $\alpha t^2+\beta t+\gamma=0$；切点是重根，从 $\gamma+\beta/2=0$ 得极线。
 \item \textbf{常规通用：}设 $x=u+rt,y=v+st$，代入一般式，令一次项为零，推出切线和极线。
\end{itemize}

\medskip
\textbf{\textcolor{CSummary}{最后三次核对}}
\begin{enumerate}[leftmargin=*,itemsep=0.25em]
 \item 标准椭圆、双曲线的中心点没有普通极线；一般式须确认 $x,y$ 系数不同时为零。
 \item $xy$、$x$、$y$ 项写成两组坐标的平均式；椭圆 $+$、双曲线 $-$。
 \item 求\textbf{切点弦}时，检查极线与曲线是否有两个不同的实交点；只有一个或没有都不能叫两切点连线。
\end{enumerate}
\end{summarybox}
